\section{Línea de tiempo de la evolución técnica de GenAI}

\subsection{2018-2021: la era del foundation model}

El mayor avance de esta etapa es el preentrenamiento + fine-tuning de Transformer. Burak menciona el GTM de esa época: primero se usan BERT y GPT-2 para resolver tareas de NLU/NLG, y después se abre el embedding a la Search empresarial y al document understanding. La práctica típica de esta etapa consiste en building an "AI co-pilot" for narrow tasks, y se enfatiza que «la preparación de datos, el feature engineering y el knowledge graph» son la verdadera ventaja competitiva.

\subsection{2022-2023: Agent y multi-turn}
El concepto de Agent comienza a popularizarse en 2022: el LLM ya no solo produce texto, sino que orchestrating tools. Tras la aparición de frameworks como Vertex AI Agent Builder y LangChain, los ingenieros pueden armar en pocos días un flujo plan→tool→feedback. El problema central de esta etapa es «hard-coded workflow + manual monitoring», por lo que Burak insiste en hacer metrics-first design.

En esta etapa es habitual que el repository incluya: prompt templates, tool manifest, evaluation harness. Burak menciona en particular que hay que documentar el sample input/output, estimated latency y failure modes de cada tool en un doc para que el planner los consulte. De lo contrario, el agent puede provocar fácilmente un cascade failure al interpretar mal un tool call.

\subsection{Después de 2024: tiempo real, control y cumplimiento}

A partir de 2024, el interés de las empresas se desplaza hacia el tiempo real y el cumplimiento normativo. El multi-modal de Gemini, el agent scheduler de Vertex y el RSP de Anthropic exigen todos un nivel más alto de «model governance». Las palabras clave de la nueva etapa son «trazabilidad», el «policy-as-code» visualizado y un «failure analysis» riguroso.

Una tendencia clave es que «legal hold + audit trail» se está convirtiendo gradualmente en baseline. El equipo de cumplimiento exige que cada decision path pueda reproducirse, e incluso que el user prompt y el tool call se encadenen en una timeline para su revisión. Muchas empresas conectan esto a un data lake y lo vinculan con el pipeline de observabilidad.

\subsection{Resumen de la sección}

Al repasar la línea de tiempo se observa que, del preentrenamiento al Agent y de ahí al cumplimiento normativo, el reto de la IA empresarial se desplaza continuamente de la capability hacia el control, lo que implica que los equipos deben iterar sin cesar sus procesos y su estructura organizacional.

